\section{Introduction}

Large language models (LLMs) excel at reasoning, coding, and multimodal tasks~\cite{olfactory2025,staticanalysis2025,egi2025}, but their hardware and deployment requirements often confine them to cloud infrastructure or high-end local systems, creating privacy, cloud-dependence, and cost concerns~\cite{energytoken2025,qwqrouting2025,tigani2025cpu,aydin2025integrated}.

Small language models (SLMs) are attractive for local deployment and lower-cost, data-sovereign operation. Fine-tuning or distillation can also make them competitive on domain-specific tasks~\cite{xie2025fusing,zhu2024distilling,rethinkllm2025}. They pair naturally with routing, verification, and planning mechanisms~\cite{aydin2025integrated,qwqrouting2025,ji2025raprag}. However, SLMs still trail frontier LLMs on broad reasoning and tasks requiring extensive contextual knowledge~\cite{egi2025,dorfner2024comparing}, which has driven recent work on hybrid and agent-based architectures that combine small and large models for a better capability--efficiency trade-off~\cite{hao2024hybrid,shao2025dot,multiagent2025}.

This is primarily a systems-design question. Across domains, hybrid and multi-agent architectures aim to improve not only accuracy but also latency, cost, and computational efficiency through edge--cloud verification, expert routing, retrieval, and planner-based orchestration~\cite{hao2024hybrid,lee2024orchestrallm,alabbasi2024teleoracle,ji2025raprag,chatzistefanidis2025symbiotic,han2025chest,energytoken2025}. These developments motivate a closer examination of how collaborative architectures allocate work between small and large models, when monolithic LLM baselines remain necessary, and when orchestration provides the best balance between capability and efficiency.

Accordingly, this review examines when hybrid or multi-agent SLM--LLM architectures provide a better trade-off between capability and efficiency than monolithic LLM baselines, and under which domain conditions routing, retrieval, planning, and role specialization allow smaller models to contribute effectively.

The main contributions are threefold: (1) a synthesis of 54 studies on hybrid SLM--LLM architectures, retrieval-centered collaboration, ensembles, and multi-agent workflows; (2) a review structure organized around four practical questions on performance, efficiency, orchestration, and domain suitability; and (3) an analytical matrix mapping each study to the research question(s) it addresses.
